% ============================================================
%  Chapter 3 — Component Description (User Documentation)
% ============================================================
\chapter{Component Documentation}
\label{ch:user-docs}

This chapter contains the documentation that is provided to end-users of
\textbf{resilient-call}.

\section{Public Class Diagram}
\label{sec:public-class-diagram}

Figure~\ref{fig:class-public} shows the public API surface — only the classes,
methods, and properties that end-users interact with directly.

\begin{figure}[H]
    \centering
    \includegraphics[width=\linewidth]{class_diagram_public}
    \caption{Public API class diagram}
    \label{fig:class-public}
\end{figure}

\section{Detailed Class and Method Descriptions}
\label{sec:api-reference}

\subsection{\texttt{CircuitBreaker}}

Wraps a callable with a three-state circuit breaker.  Thread-safe.

\begin{table}[H]
    \centering
    \caption{\texttt{CircuitBreaker} — constructor parameters}
    \begin{tabularx}{\linewidth}{>{\ttfamily}l >{\ttfamily}l X}
        \toprule
        Parameter & Type & Description \\
        \midrule
        name   & str                   & Unique name used in error messages and registry look-ups. \\
        config & CircuitBreakerConfig  & Configuration object.  Library defaults are used when omitted. \\
        \bottomrule
    \end{tabularx}
\end{table}

\begin{table}[H]
    \centering
    \caption{\texttt{CircuitBreaker} — public methods and properties}
    \begin{tabularx}{\linewidth}{>{\ttfamily}l X}
        \toprule
        Member & Description \\
        \midrule
        state              & Read-only property.  Returns the current \texttt{CircuitBreakerState}. \\
        failure\_count     & Number of consecutive failures recorded in CLOSED state. \\
        call(func, *args)  & Execute \texttt{func} with circuit protection.  Raises \texttt{CircuitBreakerError} when OPEN. \\
        reset()            & Manually force the breaker to CLOSED and clear all counters. \\
        \bottomrule
    \end{tabularx}
\end{table}

\subsection{\texttt{CircuitBreakerConfig}}

\begin{table}[H]
    \centering
    \caption{\texttt{CircuitBreakerConfig} — fields}
    \begin{tabularx}{\linewidth}{>{\ttfamily}l >{\ttfamily}c X}
        \toprule
        Field & Default & Description \\
        \midrule
        failure\_threshold   & 5           & Consecutive failures to trip OPEN. \\
        success\_threshold   & 2           & Consecutive HALF\_OPEN successes to close. \\
        timeout              & 60.0        & Seconds to remain OPEN before HALF\_OPEN. \\
        expected\_exceptions & (Exception,) & Only these exception types count as failures. \\
        on\_state\_change    & None        & Callback: \texttt{(name, old\_state, new\_state) -> None}. \\
        \bottomrule
    \end{tabularx}
\end{table}

\subsection{\texttt{RetryPolicy}}

\begin{table}[H]
    \centering
    \caption{\texttt{RetryPolicy} — fields}
    \begin{tabularx}{\linewidth}{>{\ttfamily}l >{\ttfamily}c X}
        \toprule
        Field & Default & Description \\
        \midrule
        max\_attempts        & 3              & Total calls including the first attempt. \\
        delay                & 1.0            & Base delay in seconds. \\
        backoff              & EXPONENTIAL    & Delay growth strategy. \\
        max\_delay           & 60.0           & Upper bound on computed delay. \\
        retryable\_exceptions & (Exception,)  & Only these trigger a retry. \\
        on\_retry            & None           & Callback: \texttt{(attempt, sleep, exc) -> None}. \\
        \bottomrule
    \end{tabularx}
\end{table}

\subsection{\texttt{BackoffStrategy} — Delay Formulas}

\begin{table}[H]
    \centering
    \caption{Backoff strategy delay formulas}
    \begin{tabular}{ll}
        \toprule
        Strategy & Formula \\
        \midrule
        \texttt{FIXED}       & $d = \texttt{delay}$ \\
        \texttt{EXPONENTIAL} & $d = \min(\texttt{delay} \times 2^{k-1},\ \texttt{max\_delay})$ \\
        \texttt{JITTER}      & $d = \text{uniform}(0,\ \min(\texttt{delay} \times 2^{k-1},\ \texttt{max\_delay}))$ \\
        \bottomrule
    \end{tabular}
\end{table}

where $k$ is the 1-based attempt number.

\subsection{\texttt{ResilientExecutor}}

\begin{table}[H]
    \centering
    \caption{\texttt{ResilientExecutor} — constructor parameters}
    \begin{tabularx}{\linewidth}{>{\ttfamily}l X}
        \toprule
        Parameter & Description \\
        \midrule
        retry\_policy   & \texttt{RetryPolicy} instance.  Omit for no retries. \\
        circuit\_breaker & \texttt{CircuitBreaker} instance.  Omit for no circuit breaking. \\
        fallback        & Callable receiving the raised exception; its return value is used instead of raising. \\
        \bottomrule
    \end{tabularx}
\end{table}

\section{API Documentation}
\label{sec:api-docs}

Full API documentation is generated automatically from source code annotations
using \textbf{Sphinx} with the \texttt{autodoc} and \texttt{napoleon} extensions.

To build the HTML documentation locally:

\begin{lstlisting}[language=bash, numbers=none, frame=none, backgroundcolor=\color{backcolour}]
pip install sphinx sphinx-rtd-theme sphinx-autodoc-typehints
cd docs/
sphinx-build -b html . _build/html
open _build/html/index.html
\end{lstlisting}

\section{Usage Examples}
\label{sec:usage-examples}

\subsection{Example 1 — Simple Retry Decorator}

The \texttt{@retry} decorator adds retry behaviour to any function in one line.
This example demonstrates recovering from a flaky weather API that fails
the first three attempts.

\begin{lstlisting}[language=Python, caption={Retry decorator — exponential backoff}, label={lst:retry-decorator}]
from resilient_call import retry, BackoffStrategy

@retry(
    max_attempts=4,
    delay=0.5,
    backoff=BackoffStrategy.EXPONENTIAL,
    retryable_exceptions=(ConnectionError, TimeoutError),
)
def fetch_weather(city: str) -> dict:
    return requests.get(f"https://api.weather.com/{city}",
                        timeout=5).json()

# Caller is unaware of retries — they happen transparently
weather = fetch_weather("Istanbul")
\end{lstlisting}

\subsection{Example 2 — Circuit Breaker with Fallback}

The \texttt{@circuit\_breaker} decorator trips the circuit after 3 consecutive
failures.  A fallback function provides a safe default value.

\begin{lstlisting}[language=Python, caption={Circuit breaker decorator with fallback}, label={lst:cb-decorator}]
from resilient_call import circuit_breaker

@circuit_breaker(
    "payments-api",
    failure_threshold=3,
    timeout=30.0,
    fallback=lambda exc: {"status": "unavailable", "retry": True},
)
def charge_card(amount: float, currency: str = "USD") -> dict:
    return payments_client.post("/charge",
                                json={"amount": amount,
                                      "currency": currency})
\end{lstlisting}

\subsection{Example 3 — Programmatic Executor (Retry + Circuit Breaker + Fallback)}

The \texttt{ResilientExecutor} combines both mechanisms with an explicit
fallback for maximum control.

\begin{lstlisting}[language=Python, caption={ResilientExecutor — full resilience stack}, label={lst:executor}]
from resilient_call import (
    ResilientExecutor, RetryPolicy, CircuitBreaker,
    CircuitBreakerConfig, BackoffStrategy,
)

executor = ResilientExecutor(
    retry_policy=RetryPolicy(
        max_attempts=5,
        delay=0.5,
        backoff=BackoffStrategy.JITTER,
        retryable_exceptions=(IOError, TimeoutError),
        on_retry=lambda attempt, sleep, exc:
            logger.warning(f"retry {attempt} after {sleep:.2f}s: {exc}"),
    ),
    circuit_breaker=CircuitBreaker(
        "inventory-api",
        CircuitBreakerConfig(
            failure_threshold=5,
            success_threshold=2,
            timeout=60.0,
            on_state_change=lambda name, old, new:
                alert_team(f"{name}: {old.value} -> {new.value}"),
        ),
    ),
    fallback=lambda exc: [],   # return empty list when all else fails
)

items = executor.execute(fetch_inventory, warehouse_id=42)
\end{lstlisting}

\section{User Guide}
\label{sec:user-guide}

\subsection{Choosing a Backoff Strategy}

\begin{itemize}[leftmargin=*, label=\textbullet]
    \item Use \textbf{FIXED} when the delay between retries must be
          predictable (e.g.\ SLA-driven retry windows).
    \item Use \textbf{EXPONENTIAL} for typical service-to-service calls
          where you want aggressive first retries and longer waits on
          repeated failure.
    \item Use \textbf{JITTER} when many instances of the same service
          call the same dependency — random spread avoids thundering herd.
\end{itemize}

\subsection{Choosing Failure and Success Thresholds}

A low \texttt{failure\_threshold} (e.g.\ 2--3) reacts quickly to failures but
may trip on transient blips.  A high threshold (e.g.\ 10--20) tolerates noise
but allows prolonged unhealthy calls.  A typical starting point is
$\texttt{failure\_threshold} = 5$, $\texttt{success\_threshold} = 2$,
$\texttt{timeout} = 60\text{s}$.

\subsection{Sharing Circuit Breakers Across Modules}

\begin{lstlisting}[language=Python, caption={Using the module-level registry}, label={lst:registry}]
# module_a.py
from resilient_call.registry import registry
cb = registry.get_or_create("payments-api")

# module_b.py — retrieves the same instance
from resilient_call.registry import registry
cb = registry.get("payments-api")
print(cb.state)   # shared state
\end{lstlisting}

\subsection{Idempotency Warning}

\begin{warningbox}
Retry logic assumes the wrapped operation is \textbf{idempotent} — that
executing it multiple times has the same effect as executing it once.
For non-idempotent operations (e.g.\ creating a database record, sending
an email) restrict \texttt{retryable\_exceptions} to network-layer errors
only, or ensure idempotency at the operation level (e.g.\ idempotency keys).
\end{warningbox}

\section{Installation Instructions and Dependencies}
\label{sec:installation}

\subsection{Requirements}

\begin{itemize}[leftmargin=*, label=\textbullet]
    \item Python 3.10 or higher
    \item No third-party runtime dependencies
\end{itemize}

\subsection{Installation}

\begin{lstlisting}[language=bash, numbers=none, frame=none, backgroundcolor=\color{backcolour}]
pip install resilient-call
\end{lstlisting}

To install with development and documentation extras:

\begin{lstlisting}[language=bash, numbers=none, frame=none, backgroundcolor=\color{backcolour}]
pip install "resilient-call[dev]"   # pytest, pytest-asyncio, coverage
pip install "resilient-call[docs]"  # sphinx, rtd-theme, typehints
\end{lstlisting}

\subsection{Verifying the Installation}

\begin{lstlisting}[language=Python, caption={Smoke test after installation}, label={lst:smoke-test}]
from resilient_call import __version__, CircuitBreaker
print(__version__)          # 1.0.0
cb = CircuitBreaker("test")
print(cb.state.value)       # closed
\end{lstlisting}

\section{Deployment Environment Requirements}
\label{sec:deployment-requirements}

\begin{table}[H]
    \centering
    \caption{Deployment environment requirements}
    \begin{tabular}{ll}
        \toprule
        Requirement & Detail \\
        \midrule
        Python runtime    & 3.10+ (CPython or PyPy) \\
        Operating system  & Any OS supported by Python \\
        Memory            & $< 1\,\text{MB}$ per 1\,000 registered circuit breakers \\
        External services & None \\
        \bottomrule
    \end{tabular}
\end{table}

\section{Changes Log}
\label{sec:changelog}

\subsection*{Version 1.0.0 — 2026-05-27}

\textbf{Initial release.} Features:

\begin{itemize}[leftmargin=*, label=\textbullet]
    \item \texttt{CircuitBreaker} with CLOSED / OPEN / HALF\_OPEN state machine.
    \item \texttt{RetryPolicy} with FIXED, EXPONENTIAL, and JITTER strategies.
    \item \texttt{ResilientExecutor} composing both mechanisms with fallback.
    \item \texttt{CircuitBreakerRegistry} with module-level singleton.
    \item \texttt{@retry} and \texttt{@circuit\_breaker} decorators.
    \item \texttt{asyncio} support via \texttt{execute\_async}.
    \item \texttt{on\_retry} and \texttt{on\_state\_change} callback hooks.
    \item Full Sphinx API documentation.
    \item 42 unit tests; 100\% branch coverage.
\end{itemize}

\section{License}
\label{sec:license}

\textbf{resilient-call} is released under the \textbf{MIT License}.

\begin{mdframed}
\small
\textbf{MIT License}\\[6pt]
Copyright (c) 2026 Arda Erdoğan\\[6pt]
Permission is hereby granted, free of charge, to any person obtaining a copy
of this software and associated documentation files (the ``Software''), to deal
in the Software without restriction, including without limitation the rights
to use, copy, modify, merge, publish, distribute, sublicense, and/or sell
copies of the Software, and to permit persons to whom the Software is
furnished to do so, subject to the following conditions:\\[6pt]
The above copyright notice and this permission notice shall be included in all
copies or substantial portions of the Software.\\[6pt]
THE SOFTWARE IS PROVIDED ``AS IS'', WITHOUT WARRANTY OF ANY KIND, EXPRESS OR
IMPLIED, INCLUDING BUT NOT LIMITED TO THE WARRANTIES OF MERCHANTABILITY,
FITNESS FOR A PARTICULAR PURPOSE AND NONINFRINGEMENT\@.
\end{mdframed}
